\begin{table*}[t]
\centering
\small
\setlength{\tabcolsep}{3.5pt}
\resizebox{\textwidth}{!}{%
\begin{tabular}{llllccc}
\toprule
\textbf{Method}
& \textbf{Persistent control state}
& \textbf{Unit advanced}
& \textbf{Next-step rule}
& \textbf{Adaptive}
& \textbf{Process selectable}
& \textbf{Evolving structure} \\
\midrule
Re$^3$~\cite{yang2022re3}
& Plan + draft
& Passage / revision stage
& Prescribed workflow
& \no & \no & \no \\
STORM~\cite{shao2024assisting}
& Outline + sources
& Outline section
& Outline order
& \no & \no & \no \\
CogWriter~\cite{wan2025cognitive}
& Hierarchical plan + draft
& Planned segment / revision
& Monitoring-triggered revision
& \partly & \no & \yes \\
WriteHERE~\cite{xiong2025beyond}
& Task graph
& Task node
& Adaptive task scheduling
& \yes & \no & \yes \\
IS-CoT~\cite{sun2026cot}
& Plan + draft
& Plan / Write / Reflect stage
& Recurring cycle
& \no & \no & \partly \\
\midrule
\condAgenticCogWriter\ \textbf{(Ours)}
& Draft + goals + history
& Writing process
& Monitor selects online
& \yes & \yes & \yes \\
\bottomrule
\end{tabular}%
}
\caption{Control organization in representative long-form writing systems. ``Adaptive'' indicates that the next control unit is chosen online rather than advanced by a prescribed order. ``Process selectable'' is stricter: a writing process itself is the unit that can be selected next. ``Evolving structure'' indicates an explicit plan, task, or goal structure that can change during generation. \partly\ denotes a restricted form of the property.}
\label{tab:architecture-comparison}
\end{table*}
